\chapter*{Preface}
\addcontentsline{toc}{chapter}{Preface}

This manual is a living guide to the \texttt{apfs-fastindex} research project.
Its purpose is to help a future engineer reproduce the work, understand why the
scope is narrow, and adapt the method to other projects that need filesystem
metadata faster than ordinary high-level traversal APIs can provide.

The source of truth remains the repository artifact trail:

\begin{itemize}
  \item research synthesis lives in \path{docs/research/RL-*.md}
  \item external evidence reviews live in \path{docs/research/sources/SR-*.md}
  \item controlled probes live in \path{docs/research/experiments/EX-*}
  \item implementation-facing boundaries live in \path{docs/implementation/}
  \item the current product intent lives in \path{spec.md}
\end{itemize}

This document should not promote hypotheses into design commitments. When it
summarizes a claim, it should preserve the project's claim discipline:
\emph{Spec} for public documentation, \emph{Observation} for experiment or
converging implementation evidence, and \emph{Hypothesis} for plausible but
unproven design direction.

The current manual is intentionally skeletal. It records the narrative structure
and the canonical cross-links so later research can fill chapters without
turning the repository into a second, inconsistent source of truth.
